% Proved after four complete actual-file/source audits. Qualitative anchor refinement only; no hereditary quantitative atlas or main movement.
\subsection{A generating anchor of rank equal to the geometric period}

\begin{proposition}[Proved: a rank-period generator after subdividing root nodes]
\label{D:rank-period-generating-root-anchor}
Let \(R\) be a complete excellent normal local surface algebra
over an algebraically closed coefficient and residue field \(k\)
of characteristic \(p>0\), and put
\(U=\operatorname{Spec}R-\{\mathfrak m\}\).
Let \(n\) be prime to \(p\), and let
\(\alpha\in H^2(U,\mu_n)\) be realized by a positive-rank
scalar-weight-one twisted bundle. Write \(\beta\) for its
geometric Brauer image and \(M\) for the order of \(\beta\).
Fix a regular resolution with SNC exceptional divisor
\(E=\bigcup_i E_i\). Let \(m_i\) be the order of the
generic residue of \(\beta\) on \(E_i\); order one is
allowed. At a node \(x\in E_i\cap E_j\), let \(h_x\)
be the order of its secondary residue.

Blow up every node once, writing \(Y'\to Y\) for the
result. On \(Y'\), give the strict transform of \(E_i\)
root order \(m_i\), and give the new exceptional curve
over \(x\) root order \(h_x\). Omit roots of order one.
Let \(\mathcal X'\) be this regular mixed-root stack.
There is an actual \(\mu_n\)-gerbe
\(\mathcal G'\to\mathcal X'\) restricting to the prescribed
coefficient gerbe of \(\alpha\) on \(U\), and a
scalar-weight-one locally free sheaf \(\mathcal A\) of
rank exactly \(M\) on \(\mathcal G'\) with the following
properties.

\begin{enumerate}
\item \(\mathcal A\) contains every scalar-weight-one
geometric inertia simple at every point. Thus it is a
generating bundle for the weight-one category.
\item If \(\mathcal F\) is a scalar-weight-one locally
free sheaf of rank \(r\), then \(M\mid r\).
Put \(s=r/M\), let \(q:\mathcal G'\to Y'\) be the
coarse map, and set
\[
 \mathcal B=q_*\mathcal End(\mathcal A),\qquad
 \mathcal E=q_*\mathcal Hom(\mathcal A,\mathcal F).
                                                        \tag{D.RA.1}
\]
The ordinary algebra \(\mathcal B\) is locally free of
rank \(M^2\), and \(\mathcal E\) is locally free as a
right \(\mathcal B\)-module of rank \(s\), locally for
the etale topology on \(Y'\).
For this fixed \(\mathcal A\), the corresponding Morita
functors are exact and commute with extension of
coefficient rings. Their restriction identifies
weight-one vector-bundle families of rank \(r\) with
locally free right \(\mathcal B\)-module families of
rank \(s\).
\item If the original SNC model has \(v\) components
and \(e\) nodes, the construction adds \(e\) curves
and replaces its nodes by \(2e\) nodes. It therefore
uses \(O_R(t)\) centres on the controlled SNC models
of degree-\(t\) covers. Pullbacks of the previously
chosen cycles retain their already proved cohomology
and intersection budgets.
\end{enumerate}
The construction of \(\mathcal A\) is qualitative.
It asserts no uniform bound for the cohomology,
generation, or formal pole constants of the
particular order \(\mathcal B\).
\end{proposition}

\begin{proof}
\emph{The known period-index input and the actual coefficient
gerbe.}
Theorem 2.1 of the author preprint of
\citet{C_CTOP2000LocalPeriodIndex} states that a central
division algebra of prime-to-\(p\) period \(M\) over the
fraction field of an excellent henselian two-dimensional
local domain with separably closed residue field is
cyclic of index \(M\). We use that precise local
theorem, not a function-field theorem for algebraic
surfaces. Its complete written proof is on pp.~11--12
of the retained preprint; its unramified vanishing
input is Theorem 1.3 and Corollary 1.9 there.
This period-index result is Known.

The regular puncture has injective Brauer restriction
to its function field \(K\). Hence the generic image
of \(\beta\) has the same period \(M\), and its index
is \(M\). The actual generic \(\mu_n\)-gerbe has a
weight-one twisted vector space of rank \(M\):
over a field the Kummer map has zero Picard term,
and the index is the least such positive rank.
In particular, this statement does not replace
\(\alpha\) by a chosen class with the same Brauer image.
Every weight-one vector space of rank \(r\) over this
generic gerbe has \(M\mid r\).

The residue calculation and the vanishing of the
prime-to-\(p\) Brauer group of the regular proper-local
model give
\[
           M=\operatorname{lcm}_i m_i,\qquad
           h_x\mid m_i,\ m_j.                         \tag{D.RA.2}
\]
These are precisely the assertions already proved in
\ref{D:SNC-general-Brauer-geometry}.
We next compute the new residues before using the
mixed-root extension of that result.

\emph{What one node blowup does.}
Work at a node \(x\) with regular parameters \(u,v\).
In the strictly henselian completed local ring, the
prime-to-\(p\) ramified class has the form
\[
                         c(u,v)_n,\qquad c\in\mathbb Z/n.
                                                        \tag{D.RA.3}
\]
Indeed the two-dimensional residue complex first
writes it as an unramified class plus two
unit-parameter symbols and \(c(u,v)_n\).
The units are \(n\)-th powers in this complete
strictly local ring, and its unramified Brauer class
vanishes. This is the local computation in the
complete proof of the cited period-index theorem.
The order of \(c\) is \(h_x\).

In the blowup chart \(v=uw\), the symbol becomes
\(c(u,uw)_n=c(u,w)_n+c(u,u)_n\).
The second summand is \(c(u,-1)_n=0\), since
\(-1\) is an \(n\)-th power in \(k\).
The exceptional residue is therefore the class of
\(w^{c}\) or its inverse in
\(k(w)^\times/k(w)^{\times n}\); the sign depends only
on the residue convention. Its order is exactly
\(h_x\), since the divisor of \(w\) has coefficients
\(1,-1\). The two new secondary residues again have
order \(h_x\). The old generic residues are unchanged.
Thus the only new root nodes have orders
\[
                    (m_i,h_x)\quad\hbox{or}\quad
                    (h_x,m_j),
 \qquad\hbox{commutator order }h_x.                   \tag{D.RA.4}
\]
If \(h_x=1\), the exceptional curve carries no root;
there is only the single old root inertia at each
intersection point. This case must not be treated
as a node with two old root orders.

The actual mixed-root Azumaya extension and the
Kummer correction by an extending line bundle, proved
in \ref{D:SNC-general-Brauer-geometry}, now give an
actual \(\mu_n\)-gerbe \(\mathcal G'\) restricting to
\(\alpha\). The local Azumaya construction is the
one in \citet[Proposition 3.2.1, complete pp.~10--11]
{Lieblich2011ResidualD}; only its already verified
componentwise-root application is used.
In particular the scalar coefficient band remains
\(\mu_n\), even when \(M<n\).

\emph{Extending the generic rank-\(M\) object.}
The existing algebraic generating bundle on
\(\mathcal G'\), supplied by
\ref{D:SNC-general-Brauer-geometry}, is used here only
as a qualitative coherence tool. Apply its coarse
Hom functor to the generic rank-\(M\) object.
Extend the resulting finite generic module to a
coherent module over the corresponding coherent
coarse algebra: first choose a coherent ordinary
lattice, then take the algebra span.
The exact inverse Morita functor produces a coherent
weight-one extension of the generic object.
Remove its torsion and take its double dual.

The root stack and gerbe have regular two-dimensional
etale scheme charts. A coherent reflexive module on
such a chart is locally free. Etale flat pullback
commutes with the duals, so these hulls descend,
remain scalar-weight one, and retain rank \(M\).
This gives \(\mathcal A\).
The reflexive extension and regular-surface facts
are exactly the full proofs recorded under
\citet[Tag 0AVC, Lemma 15.24.19, and Tag 0B3N,
Lemma 31.13.15]{C_StacksRootAnchorReflexive}.
This construction is made over \(k\); it is not
a claim that relative reflexive hulls preserve
flatness over arbitrary coefficient rings.

\emph{Generic multiplicities are forced.}
On the generic residual gerbe of a component of
root order \(m\), the weight-one inertia characters
form the cyclic torsor attached to its residue
character. That torsor is connected of degree \(m\),
because the residue has exact order \(m\).
Consequently a rank-\(d\) weight-one bundle has
every one of the \(m\) geometric characters with
the same multiplicity \(d/m\).
This also follows directly by Galois permutation
of the character eigenspaces. The exact residual
category used here is
\ref{D:SNC-general-residual-category}.
Multiplicities are locally constant along the
whole component, including its smooth point on
each new node branch.

\emph{At a subdivided node the marginal determines
every simple.}
We give the mixed-order algebra, to avoid inferring
a generating property just from total rank.
At a geometric node of type \((m,h)\) in
(D.RA.4), the central extension of
\(\mu_m\times\mu_h\) by the scalar band has
scalar-weight-one commutator of exact order \(h\).
After scalar normalizations its twisted group
algebra is generated by \(X,Y\) with
\[
        X^m=1,\qquad Y^h=1,\qquad
        XY=\zeta YX,\quad \operatorname{ord}(\zeta)=h.
                                                        \tag{D.RA.5}
\]
The scalar normalizations are made over the
algebraically closed geometric residue field.
They do not assert a splitting of the central
inertia extension itself.

The \(m\) eigencharacters of \(X\) break into
disjoint orbits of size \(h\) under \(Y\).
For each orbit there is exactly one irreducible
module, of dimension \(h\): choose a vector in
one eigenspace, apply \(Y\), and normalize its
cyclic product using \(Y^h=1\).
Different orbits give nonisomorphic modules,
and all modules are semisimple since the inertia
order is prime to \(p\).
Thus there are \(m/h\) simples. Each restricts
to exactly one orbit, once at each of its
\(h\) branch characters. This is the explicit
mixed-order special case of the node algebra
already proved in
\ref{D:SNC-general-residual-category}; compare the
central-extension and commutator computations in
\citet[Lemma 4.3.1 and Proposition 4.3.2, complete
pp.~13--14]{Lieblich2011ResidualD}.

For \(\mathcal A\), every branch character has
multiplicity \(M/m\). Equation (D.RA.5) therefore
forces every node simple to occur with multiplicity
\(M/m>0\). For \(\mathcal F\), the same multiplicity
is \(r/m=s(M/m)\). The other new node type is
identical after interchanging the two branches.
At a point with only one root inertia the same
conclusion is the generic character calculation;
at a point without root inertia it is immediate.
This proves the generating assertion and, more
precisely, that every geometric inertia fibre of
\(\mathcal F\) is isomorphic to that of
\(\mathcal A^{\oplus s}\).

\emph{Fixed-anchor Morita and coefficient bases.}
The preceding inertia assertion can be checked on
strictly local root charts. There the finite tame
inertia representation splits into its character
or simple factors, and taking invariants is exact.
An isomorphism of the two residual representations
lifts to an equivariant map
\(\mathcal A^{\oplus s}\to\mathcal F\);
Nakayama and equal rank make it an isomorphism
near that geometric point. The same argument
applies to coefficient algebras: the fixed finite
inertia has invertible order, so invariants and
the invariant Hom commute with coefficient
extension, and the geometric fibre isomorphism
lifts on the corresponding local frame opens.
No saturation of a varying module is involved.

The scalar centre cancels in
\(\mathcal End(\mathcal A)\); its remaining
equivariance factors through the abelian root
quotient. A homogeneous basis, adjusted by the
unique root monomial in each residue class,
shows that its coarse invariant module is free
of rank \(M^2\) on each local root chart.
This is the same all-ring monomial calculation
as in the existing mixed-root presentation.
The local isomorphism
\(\mathcal F\simeq\mathcal A^{\oplus s}\)
therefore identifies
\(\mathcal E\simeq\mathcal B^{\oplus s}\).

Equivalently, on these charts \(\mathcal A\) is
a projective generator for the weight-one
equivariant category: its Hom is exact and it
contains all residual simples. The usual
evaluation and coevaluation maps give the
Morita equivalence. The maps and their inverse
identities commute with restriction and hence
descend. All the operations are fixed finite
module maps, invariant direct summands, and
tensor products, so the equivalence commutes
with coefficient extension. In the locally
free subcategories its inverse is visibly
locally \(\mathcal A^{\oplus s}\).
This establishes the precise family assertion
in item~2 for the already fixed anchor.

Finally, the original SNC nodes are pairwise
distinct and have no triple points. One blowup
adds one smooth rational component and replaces
that node by two transverse nodes. Old strict
transforms are mutually disjoint, as are new
exceptional curves. Thus the stated centre
and graph counts are exact. The already proved
pullback comparison for the inherited cycles
applies to these ordinary point blowups.
\end{proof}

\begin{remark}[The quantitative question that remains]
The index theorem supplies a generic cyclic splitting
extension \(K(f^{1/M})/K\), but its proof chooses
\(\operatorname{div}(f)\) equal to the ramification
divisor plus an auxiliary divisor avoiding finitely
many points. It does not control that auxiliary
divisor's degrees, and that extension need not be
etale on \(U\). No bounded-geometry splitting cover
is inferred here.

More importantly, \(\mathcal B=q_*\mathcal End(\mathcal A)\)
in (D.RA.1) need not be the canonical full pullback
of its punctured algebra. The same distinction
applies to \(\mathcal E\).
The full-bundle degree and generation theorems
for the canonical ordinary pullbacks do not
automatically transfer to these particular
coarse lattices. Even the proved bounded
divisorial sandwiches do not imply generation
of an intermediate lattice or flatness of a
relative saturation.

A possible whole-cycle presentation would require,
in addition, controlled \(\mathcal B\)-generating
quotients of \(\mathcal E\), short exact syzygies
in that locally free module category, and fixed
Hom complexes with total size \(O_R(tr^2)\).
The local rank identity
\(M^2(r/M)^2=r^2\) is consistent with that target,
but is not a proof of any of those global
requirements. This proposition is an auxiliary
qualitative refinement, not removal of the
unbounded-period hereditary counting gate.
\end{remark}
